\IfFileExists{stacks-project.cls}{%
\documentclass{stacks-project}
}{%
\documentclass{amsart}
}

% For dealing with references we use the comment environment
\usepackage{verbatim}
\newenvironment{reference}{\comment}{\endcomment}
%\newenvironment{reference}{}{}
\newenvironment{slogan}{\comment}{\endcomment}
\newenvironment{history}{\comment}{\endcomment}

% For commutative diagrams we use Xy-pic
\usepackage[all]{xy}

% We use 2cell for 2-commutative diagrams.
\xyoption{2cell}
\UseAllTwocells

% We use multicol for the list of chapters between chapters
\usepackage{multicol}

% This is generally recommended for better output
\usepackage{lmodern}
\usepackage[T1]{fontenc}

% For cross-file-references
\usepackage{xr-hyper}

% Package for hypertext links:
\usepackage{hyperref}

% For any local file, say "hello.tex" you want to link to please
% use \externaldocument[hello-]{hello}
\externaldocument[introduction-]{introduction}
\externaldocument[conventions-]{conventions}
\externaldocument[sets-]{sets}
\externaldocument[categories-]{categories}
\externaldocument[topology-]{topology}
\externaldocument[sheaves-]{sheaves}
\externaldocument[sites-]{sites}
\externaldocument[stacks-]{stacks}
\externaldocument[fields-]{fields}
\externaldocument[algebra-]{algebra}
\externaldocument[brauer-]{brauer}
\externaldocument[homology-]{homology}
\externaldocument[derived-]{derived}
\externaldocument[simplicial-]{simplicial}
\externaldocument[more-algebra-]{more-algebra}
\externaldocument[smoothing-]{smoothing}
\externaldocument[modules-]{modules}
\externaldocument[sites-modules-]{sites-modules}
\externaldocument[injectives-]{injectives}
\externaldocument[cohomology-]{cohomology}
\externaldocument[sites-cohomology-]{sites-cohomology}
\externaldocument[dga-]{dga}
\externaldocument[dpa-]{dpa}
\externaldocument[sdga-]{sdga}
\externaldocument[hypercovering-]{hypercovering}
\externaldocument[schemes-]{schemes}
\externaldocument[constructions-]{constructions}
\externaldocument[properties-]{properties}
\externaldocument[morphisms-]{morphisms}
\externaldocument[coherent-]{coherent}
\externaldocument[divisors-]{divisors}
\externaldocument[limits-]{limits}
\externaldocument[varieties-]{varieties}
\externaldocument[topologies-]{topologies}
\externaldocument[descent-]{descent}
\externaldocument[perfect-]{perfect}
\externaldocument[more-morphisms-]{more-morphisms}
\externaldocument[flat-]{flat}
\externaldocument[groupoids-]{groupoids}
\externaldocument[more-groupoids-]{more-groupoids}
\externaldocument[etale-]{etale}
\externaldocument[chow-]{chow}
\externaldocument[intersection-]{intersection}
\externaldocument[pic-]{pic}
\externaldocument[weil-]{weil}
\externaldocument[adequate-]{adequate}
\externaldocument[dualizing-]{dualizing}
\externaldocument[duality-]{duality}
\externaldocument[discriminant-]{discriminant}
\externaldocument[derham-]{derham}
\externaldocument[local-cohomology-]{local-cohomology}
\externaldocument[algebraization-]{algebraization}
\externaldocument[curves-]{curves}
\externaldocument[resolve-]{resolve}
\externaldocument[models-]{models}
\externaldocument[functors-]{functors}
\externaldocument[equiv-]{equiv}
\externaldocument[pione-]{pione}
\externaldocument[etale-cohomology-]{etale-cohomology}
\externaldocument[proetale-]{proetale}
\externaldocument[relative-cycles-]{relative-cycles}
\externaldocument[more-etale-]{more-etale}
\externaldocument[trace-]{trace}
\externaldocument[crystalline-]{crystalline}
\externaldocument[spaces-]{spaces}
\externaldocument[spaces-properties-]{spaces-properties}
\externaldocument[spaces-morphisms-]{spaces-morphisms}
\externaldocument[decent-spaces-]{decent-spaces}
\externaldocument[spaces-cohomology-]{spaces-cohomology}
\externaldocument[spaces-limits-]{spaces-limits}
\externaldocument[spaces-divisors-]{spaces-divisors}
\externaldocument[spaces-over-fields-]{spaces-over-fields}
\externaldocument[spaces-topologies-]{spaces-topologies}
\externaldocument[spaces-descent-]{spaces-descent}
\externaldocument[spaces-perfect-]{spaces-perfect}
\externaldocument[spaces-more-morphisms-]{spaces-more-morphisms}
\externaldocument[spaces-flat-]{spaces-flat}
\externaldocument[spaces-groupoids-]{spaces-groupoids}
\externaldocument[spaces-more-groupoids-]{spaces-more-groupoids}
\externaldocument[bootstrap-]{bootstrap}
\externaldocument[spaces-pushouts-]{spaces-pushouts}
\externaldocument[spaces-chow-]{spaces-chow}
\externaldocument[groupoids-quotients-]{groupoids-quotients}
\externaldocument[spaces-more-cohomology-]{spaces-more-cohomology}
\externaldocument[spaces-simplicial-]{spaces-simplicial}
\externaldocument[spaces-duality-]{spaces-duality}
\externaldocument[formal-spaces-]{formal-spaces}
\externaldocument[restricted-]{restricted}
\externaldocument[spaces-resolve-]{spaces-resolve}
\externaldocument[formal-defos-]{formal-defos}
\externaldocument[defos-]{defos}
\externaldocument[cotangent-]{cotangent}
\externaldocument[examples-defos-]{examples-defos}
\externaldocument[algebraic-]{algebraic}
\externaldocument[examples-stacks-]{examples-stacks}
\externaldocument[stacks-sheaves-]{stacks-sheaves}
\externaldocument[criteria-]{criteria}
\externaldocument[artin-]{artin}
\externaldocument[quot-]{quot}
\externaldocument[stacks-properties-]{stacks-properties}
\externaldocument[stacks-morphisms-]{stacks-morphisms}
\externaldocument[stacks-limits-]{stacks-limits}
\externaldocument[stacks-cohomology-]{stacks-cohomology}
\externaldocument[stacks-perfect-]{stacks-perfect}
\externaldocument[stacks-introduction-]{stacks-introduction}
\externaldocument[stacks-more-morphisms-]{stacks-more-morphisms}
\externaldocument[stacks-geometry-]{stacks-geometry}
\externaldocument[moduli-]{moduli}
\externaldocument[moduli-curves-]{moduli-curves}
\externaldocument[examples-]{examples}
\externaldocument[exercises-]{exercises}
\externaldocument[guide-]{guide}
\externaldocument[desirables-]{desirables}
\externaldocument[coding-]{coding}
\externaldocument[obsolete-]{obsolete}
\externaldocument[fdl-]{fdl}
\externaldocument[index-]{index}

% Theorem environments.
%
\theoremstyle{plain}
\newtheorem{theorem}[subsection]{Theorem}
\newtheorem{proposition}[subsection]{Proposition}
\newtheorem{lemma}[subsection]{Lemma}

\theoremstyle{definition}
\newtheorem{definition}[subsection]{Definition}
\newtheorem{example}[subsection]{Example}
\newtheorem{exercise}[subsection]{Exercise}
\newtheorem{situation}[subsection]{Situation}

\theoremstyle{remark}
\newtheorem{remark}[subsection]{Remark}
\newtheorem{remarks}[subsection]{Remarks}

\numberwithin{equation}{subsection}

% Macros
%
\def\lim{\mathop{\mathrm{lim}}\nolimits}
\def\colim{\mathop{\mathrm{colim}}\nolimits}
\def\Spec{\mathop{\mathrm{Spec}}}
\def\Hom{\mathop{\mathrm{Hom}}\nolimits}
\def\Ext{\mathop{\mathrm{Ext}}\nolimits}
\def\SheafHom{\mathop{\mathcal{H}\!\mathit{om}}\nolimits}
\def\SheafExt{\mathop{\mathcal{E}\!\mathit{xt}}\nolimits}
\def\Sch{\mathit{Sch}}
\def\Mor{\mathop{\mathrm{Mor}}\nolimits}
\def\Ob{\mathop{\mathrm{Ob}}\nolimits}
\def\Sh{\mathop{\mathit{Sh}}\nolimits}
\def\NL{\mathop{N\!L}\nolimits}
\def\CH{\mathop{\mathrm{CH}}\nolimits}
\def\proetale{{pro\text{-}\acute{e}tale}}
\def\etale{{\acute{e}tale}}
\def\QCoh{\mathit{QCoh}}
\def\Ker{\mathop{\mathrm{Ker}}}
\def\Im{\mathop{\mathrm{Im}}}
\def\Coker{\mathop{\mathrm{Coker}}}
\def\Coim{\mathop{\mathrm{Coim}}}

% Boxtimes
%
\DeclareMathSymbol{\boxtimes}{\mathbin}{AMSa}{"02}

%
% Macros for moduli stacks/spaces
%
\def\QCohstack{\mathcal{QC}\!\mathit{oh}}
\def\Cohstack{\mathcal{C}\!\mathit{oh}}
\def\Spacesstack{\mathcal{S}\!\mathit{paces}}
\def\Quotfunctor{\mathrm{Quot}}
\def\Hilbfunctor{\mathrm{Hilb}}
\def\Curvesstack{\mathcal{C}\!\mathit{urves}}
\def\Polarizedstack{\mathcal{P}\!\mathit{olarized}}
\def\Complexesstack{\mathcal{C}\!\mathit{omplexes}}
% \Pic is the operator that assigns to X its picard group, usage \Pic(X)
% \Picardstack_{X/B} denotes the Picard stack of X over B
% \Picardfunctor_{X/B} denotes the Picard functor of X over B
\def\Pic{\mathop{\mathrm{Pic}}\nolimits}
\def\Picardstack{\mathcal{P}\!\mathit{ic}}
\def\Picardfunctor{\mathrm{Pic}}
\def\Deformationcategory{\mathcal{D}\!\mathit{ef}}
\usepackage{cleveref}
\crefname{originallemma}{original lemma}{original lemmas}
\Crefname{originallemma}{Original lemma}{Original lemmas}
\makeatletter
\newlabel{lemma-f-shriek-composition}{{0F50}{source}{}{sourcepointer-1}{}}
\newlabel{lemma-f-shriek-composition@cref}{{[originallemma][0][]0F50}{[1][0][]source}{}{}{}}
\makeatother
\crefname{originallemma}{original lemma}{original lemmas}
\Crefname{originallemma}{Original lemma}{Original lemmas}
\makeatletter
\newlabel{lemma-lqf-base-change-f-shriek}{{0F5J}{source}{}{sourcepointer-2}{}}
\newlabel{lemma-lqf-base-change-f-shriek@cref}{{[originallemma][0][]0F5J}{[1][0][]source}{}{}{}}
\makeatother
\crefname{originallemma}{original lemma}{original lemmas}
\Crefname{originallemma}{Original lemma}{Original lemmas}
\makeatletter
\newlabel{lemma-lqf-colimit-f-shriek}{{0F5H}{source}{}{sourcepointer-3}{}}
\newlabel{lemma-lqf-colimit-f-shriek@cref}{{[originallemma][0][]0F5H}{[1][0][]source}{}{}{}}
\makeatother
\crefname{originallemma}{original lemma}{original lemmas}
\Crefname{originallemma}{Original lemma}{Original lemmas}
\makeatletter
\newlabel{lemma-lqf-f-shriek-stalk}{{0F5F}{source}{}{sourcepointer-4}{}}
\newlabel{lemma-lqf-f-shriek-stalk@cref}{{[originallemma][0][]0F5F}{[1][0][]source}{}{}{}}
\makeatother
\crefname{originallemma}{original lemma}{original lemmas}
\Crefname{originallemma}{Original lemma}{Original lemmas}
\makeatletter
\newlabel{lemma-lqf-separated-shriek-composition}{{0F79}{source}{}{sourcepointer-5}{}}
\newlabel{lemma-lqf-separated-shriek-composition@cref}{{[originallemma][0][]0F79}{[1][0][]source}{}{}{}}
\makeatother
\title{Composition isomorphism (gf)\_! equals g\_! f\_! for locally quasi-finite morphisms with five compatibility properties}
\author{The Stacks Project Authors}
\date{stacks-project:more-etale.tex:local-snapshot-20261012; GNU FDL 1.2 or later}
\begin{document}
\maketitle
\setcounter{subsection}{1}
\noindent\textit{Editorial note.} All delivered mathematical text is the author's own, copied from the pinned Stacks Project chapter file more-etale.tex; no statement or proof step was rewritten, reordered or removed. The counted span is the author environment and its own proof environment. References inside the span are handed to the assembler as navigation stubs that carry the canonical tag of the cited tag only; whether each of them is load-bearing is recorded on the stub itself.\par\medskip
\noindent\textit{Editorial note.} Disclosure: two soft spots in the property checks. Property (3): after the fibre-wise identification with the separated-case map the author writes 'The reader then immediately sees that we obtain property (3)' (proof line 1963) -- the square is displayed and the horizontal maps are named. Property (5): 'This is clear from the corresponding compatibility of the maps constructed in Lemma lemma-lqf-separated-shriek-composition' (proof lines 1954-1955). Both rely on displayed data and named lemmas; no verbatim 'we omit' occurs in the proof.\par\medskip
\noindent\textit{Allowed standard prerequisites (editorial).}\begin{itemize}
\item Colimit exactness for locally quasi-finite lower shriek (Lemma lemma-lqf-colimit-f-shriek, canonical tag 0F5H). cited at proof line 1922. Gives exactness of the top and bottom rows of the construction diagram.
\item Separated composition isomorphism (Lemma lemma-lqf-separated-shriek-composition, canonical tag 0F79). cited at proof lines 1893-1895, 1923, 1946 and 1949. Provides the separated case, the commuting left squares, and the properties (1)(2)(5) compatibilities.
\end{itemize}\medskip
\ifdefined\hypertarget\hypertarget{sourcepointer-1}{}\fi
\noindent\textit{Source reference.} Original lemma 0F50: sources/stacks-chow/tags.txt, source lines 18920--18920. Cited at proof lines 1892, 1950 and 1962: separated sub-case, property (1) compatibility, and the fibre-wise identification.\par\medskip
\ifdefined\hypertarget\hypertarget{sourcepointer-2}{}\fi
\noindent\textit{Source reference.} Original lemma 0F5J: sources/stacks-chow/tags.txt, source lines 18939--18939. Cited at statement line 1886: the horizontal maps of the property (5) diagram.\par\medskip
\ifdefined\hypertarget\hypertarget{sourcepointer-3}{}\fi
\noindent\textit{Source reference.} Original lemma 0F5H: sources/stacks-chow/tags.txt, source lines 18937--18937. Cited at proof line 1922: exactness of the top and bottom rows of the construction diagram.\par\medskip
\ifdefined\hypertarget\hypertarget{sourcepointer-4}{}\fi
\noindent\textit{Source reference.} Original lemma 0F5F: sources/stacks-chow/tags.txt, source lines 18935--18935. Cited at statement line 1854: the horizontal arrows of the property (3) diagram.\par\medskip
\ifdefined\hypertarget\hypertarget{sourcepointer-5}{}\fi
\noindent\textit{Source reference.} Original lemma 0F79: sources/stacks-chow/tags.txt, source lines 18999--18999. Cited at proof lines 1894, 1923, 1946, 1949 and 1955: the separated case and its compatibilities used throughout.\par\medskip
\begin{lemma}
\label{lemma-lqf-shriek-composition}
Let $f : X \to Y$ and $g : Y \to Z$ be composable locally quasi-finite
morphisms of schemes. Then there is a canonical isomorphism of functors
$$
(g \circ f)_! \longrightarrow g_! \circ f_!
$$
These isomorphisms satisfy the following properties:
\begin{enumerate}
\item If $f$ and $g$ are separated, then the isomorphism agrees
with Lemma \ref{lemma-f-shriek-composition}.
\item If $g$ is separated, then the isomorphism agrees with
Lemma \ref{lemma-lqf-separated-shriek-composition}.
\item For a geometric point $\overline{z} : \Spec(k) \to Z$ the diagram
$$
\xymatrix{
((g \circ f)_!\mathcal{F})_{\overline{z}} \ar[d] \ar[rr] & &
\bigoplus\nolimits_{g(f(\overline{x})) = \overline{z}}
\mathcal{F}_{\overline{x}} \ar@{=}[d] \\
(g_!f_!\mathcal{F})_{\overline{z}} \ar[r] &
\bigoplus\nolimits_{g(\overline{y}) = \overline{z}}
(f_!\mathcal{F})_{\overline{y}} \ar[r] &
\bigoplus\nolimits_{g(f(\overline{x})) = \overline{z}}
\mathcal{F}_{\overline{x}}
}
$$
is commutative where the horizontal arrows are given by
Lemma \ref{lemma-lqf-f-shriek-stalk}.
\item Let $h : Z \to T$ be a third locally quasi-finite
morphism of schemes. Then the diagram
$$
\xymatrix{
(h \circ g \circ f)_! \ar[r] \ar[d] &
(h \circ g)_! \circ f_! \ar[d] \\
h_! \circ (g \circ f)_! \ar[r] &
h_! \circ g_! \circ f_!
}
$$
commutes.
\item Suppose that we have a diagram of schemes
$$
\xymatrix{
X' \ar[d]_{f'} \ar[r]_c & X \ar[d]^f \\
Y' \ar[d]_{g'} \ar[r]_b & Y \ar[d]^g \\
Z' \ar[r]^a & Z
}
$$
with both squares cartesian and $f$ and $g$
locally quasi-finite. Then the diagram
$$
\xymatrix{
a^{-1} \circ (g \circ f)_! \ar[d] \ar[rr] & &
(g' \circ f')_! \circ c^{-1} \ar[d] \\
a^{-1} \circ g_! \circ f_! \ar[r] &
g'_! \circ b^{-1} \circ f_! \ar[r] &
g'_! \circ f'_! \circ c^{-1}
}
$$
commutes where the horizontal arrows are those of
Lemma \ref{lemma-lqf-base-change-f-shriek}.
\end{enumerate}
\end{lemma}

\begin{proof}
If $f$ and $g$ are separated, then this is a special case of
Lemma \ref{lemma-f-shriek-composition}.
If $g$ is separated, then this is a special case of
Lemma \ref{lemma-lqf-separated-shriek-composition}
which moreover agrees with the case where $f$ and $g$ are separated.

\medskip\noindent
Construction in the general case. Choose an open covering $Y = \bigcup Y_i$
such that the restriction $g_i : Y_i \to Z$ of $g$ is separated.
Set $X_i = f^{-1}(Y_i)$ and denote $f_i : X_i \to Y_i$ the restriction
of $f$. Also denote $h = g \circ f$ and $h_i : X_i \to Z$ the restriction
of $h$. Consider the following diagram
$$
\xymatrix{
\bigoplus\nolimits_{i_0, i_1}
h_{i_0i_1, !}\mathcal{F}|_{X_{i_0i_1}} \ar[r] \ar[d] &
\bigoplus\nolimits_{i_0} h_{i_0, !}\mathcal{F}|_{X_{i_0}} \ar[r] \ar[d] &
h_!\mathcal{F} \ar[r] \ar@{..>}[dd] &
0 \\
\bigoplus\nolimits_{i_0, i_1}
g_{i_0i_1, !} f_{i_0i_1, !}\mathcal{F}|_{X_{i_0i_1}} \ar[r] \ar[d] &
\bigoplus\nolimits_{i_0}
g_{i_0, !} f_{i_0, !}\mathcal{F}|_{X_{i_0}} \ar[d] \\
\bigoplus\nolimits_{i_0, i_1}
g_{i_0i_1, !} (f_!\mathcal{F})|_{Y_{i_0i_1}} \ar[r] &
\bigoplus\nolimits_{i_0}
g_{i_0, !} (f_!\mathcal{F})|_{Y_{i_0}} \ar[r] &
g_!f_!\mathcal{F} \ar[r] &
0
}
$$
By Lemma \ref{lemma-lqf-colimit-f-shriek} the top and bottom row
in the diagram are exact. By Lemma \ref{lemma-lqf-separated-shriek-composition}
the top left square commutes. The vertical arrows in the
lower left square come about because
$(f_!\mathcal{F})|_{Y_{i_0i_1}} = f_{i_0i_1, !}\mathcal{F}|_{X_{i_0i_1}}$ and
$(f_!\mathcal{F})|_{Y_{i_0}} = f_{i_0, !}\mathcal{F}|_{X_{i_0}}$
as the construction of $f_!$ is local on the base. Moreover, these
equalities are (of course) compatible with the identifications
$((f_!\mathcal{F})|_{Y_{i_0}})|_{Y_{i_0i_1}} =
(f_!\mathcal{F})|_{Y_{i_0i_1}}$ and
$(f_{i_0, !}\mathcal{F}|_{X_{i_0}})|_{Y_{i_0i_1}} =
f_{i_0i_1, !}\mathcal{F}|_{X_{i_0i_1}}$
which are used (together with the covariance for open embeddings
for $Y_{i_0i_1} \subset Y_{i_0}$)
to define the horizontal maps of the lower left square.
Thus this square commutes as well.
In this way we conclude there is a unique
dotted arrow as indicated in the diagram and
moreover this arrow is an isomorphism.

\medskip\noindent
Proof of properties (1) -- (5). Fix the open covering $Y = \bigcup Y_i$.
Observe that if $Y \to Z$ happens to be separated, then we get a dotted
arrow fitting into the huge diagram above by using the map of
Lemma \ref{lemma-lqf-separated-shriek-composition}
(by the very properties of that lemma).
This proves (2) and hence also (1) by the compatibility of the
maps of Lemma \ref{lemma-lqf-separated-shriek-composition}
and Lemma \ref{lemma-f-shriek-composition}.
Next, for any scheme $Z'$ over $Z$, we obtain the compatibility in (5)
for the map $(g' \circ f')_! \to g'_! \circ f'_!$
constructed using the open covering $Y' = \bigcup b^{-1}(Y_i)$.
This is clear from the corresponding compatibility of the maps
constructed in Lemma \ref{lemma-lqf-separated-shriek-composition}.
In particular, we can consider a geometric point
$\overline{z} : \Spec(k) \to Z$. Since
$X_{\overline{z}} \to Y_{\overline{z}} \to \Spec(k)$
are separated maps, we find that the base change of
$(g \circ f)_!\mathcal{F} \to g_! f_! \mathcal{F}$
by $\overline{z}$ is equal to the map of
Lemma \ref{lemma-f-shriek-composition}.
The reader then immediately sees that we obtain property (3).
Of course, property (3) guarantees that our transformation of functors
$(g \circ f)_! \to g_! \circ f_!$ constructed using the open covering
$Y = \bigcup Y_i$ doesn't depend on the choice of this open covering.
Finally, property (4) follows by looking at what happens on stalks
using the already proven property (3).
\end{proof}

\end{document}
